\documentclass[../../main.tex]{subfiles}

\begin{document}

\subsection{Obiettivi scientifici} \label{subsec:obiettivi-scientifici}
  Il programma scientifico di DUNE è diviso in obiettivi primari: quelli per
  i quali l'esperimento è effettivamente realizzato; e secondari: obiettivi
  raggiungibili dall'esperimento in maniera ``accidentale''.
  I primi sono~\cite{dunecollaboration2022snowmassneutrinofrontierdune}:
  \begin{itemize}
    \item effettuare misure precise, utilizzando fasci di $\Pnum$ e $\APnum$,
    dei parametri di oscillazione tra $\nu_1 - \nu_3$ e $\nu_2 - \nu_3$, inclusa
    la fase di violazione di CP $\delta_{CP}$, determinare l'ordinamento delle
    masse dei neutrini e testare il paradigma a tre neutrini;

    \item studiare i neutrini emessi da supernove a collasso nucleare
    (\textit{core-collapse}), che potrebbero spiegare l'evoluzione
    dell'universo;

    \item cercare il decadimento del protone, previsto dalla
    maggior parte delle Teorie di Grande Unificazione (GUT).
  \end{itemize}
  Gli obiettivi secondari sono invece:
  \begin{itemize}
    \item effettuare ulteriori misure di interazioni oltre il Modello Standard;

    \item misurare le oscillazioni dei neutrini atmosferici;

    \item effettuare misure di altri fenomeni astrofisici.
  \end{itemize}


\end{document}